\begin{enumerate}
%1
\item \textbf{อนุพันธ์ของฟังก์ชันประกอบ}

จงเขียนฟังก์ชันต่อไปนี้ ให้เป็นฟังก์ชันประกอบ

\begin{tasks}[after-skip = 80pt](3)
\task $f(x) = (2x-1)^4$
\task $f(x) = \frac{1}{(x^2+1)^5}$
\task $f(x) = \sqrt{x^2+2x+1}$
\end{tasks}

ถ้ากำหนดให้ $f(x) = (u \circ v)(x)$ แล้วอนุพันธ์ของ $f(x)$ หาได้ด้วยกฎลูกโซ่ดังนี้

\[ \dx f(x) = \dx (u \circ v)(x) = \blank\]

จงหาอนุพันธ์ของฟังก์ชันต่อไปนี้
\begin{tasks}[after-skip = 80pt](3)
\task $f(x) = (3x+2)^5$
\task $f(x) = \frac{1}{(x^4+x)^3}$
\task $f(x) = \sqrt[4]{x^3 + 1}$
\end{tasks}

เราสามารถสรุปกฏลูกโซ่สำหรับฟังก์ชันเลขยกกำลังได้ดังนี้

\[ \dx (u(x))^n =  \blank\]

%2
\item \textbf{อนุพันธ์ของฟังก์ชันอดิศัย}

จงหาอนุพันธ์ของฟังก์ชันต่อไปนี้
\begin{tasks}[after-skip=80pt, after-item-skip=80pt](3)
\task $f(x) = 2e^{3x}$
\task $f(x) = e^{\frac{1}{x}}$
\task $f(x) = e^{\sqrt{x^2+1}}$
\task $f(x) = 4\ln( 5x+6)$
\task $f(x) = \ln \left( \frac{1}{x} \right)$
\task $f(x) = \ln (e^x)$
\end{tasks}

\newpage

\item \textbf{อนุพันธ์ของฟังก์ชังตรีโกณมิติ}

\[ \dx \sin (x) = \blank \qquad \dx \cos(x) = \blank\]

เมื่อใช้ร่วมกับกฏลูกโซ่แล้ว จะได้
\[ \dx \sin (u(x)) = \blank \qquad \dx \cos (u(x)) = \blank\]

จงหาอนุพันธ์ของฟังก์ชันตรีโกณมิติต่อไปนี้
\begin{tasks}[after-skip=80pt, after-item-skip=80pt](3)
\task $f(x) = \sin(4x)$
\task $f(x) = \cos(-3x)$
\task $f(x) = \sin(x^3)$
\task $f(x) = \cos\left(\frac{1}{x}\right)$
\task $f(x) = \sin(4x-3) + 2$
\task $f(x) = \cos(x^3+x^2+x+1)$
\task $f(x) = \sin(3x) + \cos (3x) $
\task $f(x) = x^2 + \sin(x^2)$
\task $f(x) = e^{2x} - \cos(2x)$
\end{tasks}

\item \textbf{โจทย์ระคน}

จงหาอนุพันธ์ของฟังก์ชันต่อไปนี้ โดยใช้สูตรที่เหมาะสม
\begin{tasks}[after-skip=120pt, after-item-skip=120pt](3)
\task $f(x) = (\sin(x))^2)$ 
\task $f(x) = \sin(e^x)$
\task $f(x) = (e^x)^3$
\task $f(x) = (\cos x)(\sin x)$
\task $f(x) = e^x (\cos x)$
\task $f(x) = \frac{\sin x}{e^x}$
\end{tasks}

\end{enumerate}

\end{document}